\subsection{\texorpdfstring{Vectorial measures with base \(\mu\)}{Vectorial measures with base mu}}

DEFINITION 4. --- Let \(\mu\) be a positive measure on T. A vectorial measure m on T, with values in F, is said to be a measure with base \(\mu\) if there exists a mapping f of T into F, scalarly locally \(\mu\) -integrable, such that \(\mathbf{m}(g) = \int g\mathbf{f}d\mu\) for every function \(g \in \mathcal{K}(\mathrm{T})\) . One then says that f is a density of m with respect to \(\mu\) , and one writes \(\mathbf{m} = \mathbf{f} \cdot \mu\) .

It is immediate that if \(f_{1}\) and \(f_{2}\) are two densities of m with respect to \(\mu\) , then \(f_{1}-f_{2}\) is scalarly locally \(\mu\) -negligible (Ch. V, §5, No.~3, Cor. 2 of Prop. 3); recall that in general this does not imply that \(f_{1}-f_{2}\) is zero locally almost everywhere (cf.~§1, Exer. 12 and No.~1, Remark 2).

Let \(\mathbf{m}\) be a measure with base \(\mu\) , of density \(\mathbf{f}\) . For a numerical function \(g\) to be essentially integrable for \(\mathbf{m}\) , it is necessary and sufficient that \(gf\) be scalarly essentially \(\mu\) -integrable (Ch. V, §5, No.~3, Th. 1).

PROPOSITION 7. --- Let f be a mapping scalarly locally integrable with respect to a positive measure \(\mu\) on T, such that, for every function \(g \in \mathcal{K}(T)\) , one has \(\int g\mathbf{f} d\mu \in F\) . Then the mapping \(g \mapsto \int g\mathbf{f} d\mu\) of \(\mathcal{K}(T)\) into F is a vectorial measure on T, with base \(\mu\) and density f with respect to \(\mu\) .

For (No.~1, Remark 3), it suffices to show that, setting \(\mathbf{m}(g) = \int g\mathbf{f}d\mu\) , \(\mathbf{z}' \circ \mathbf{m}\) is a scalar measure for every \(\mathbf{z}' \in \mathrm{F}'\) . But since \(\mathbf{z}'(\mathbf{m}(g)) = \int g\langle \mathbf{z}', \mathbf{f} \rangle d\mu\) , one has \(\mathbf{z}' \circ \mathbf{m} = \langle \mathbf{z}', \mathbf{f} \rangle \cdot \mu\) , whence our assertion.

PROPOSITION 8. --- Let \(\mu\) be a positive measure on T, m a measure on T with values in F, with base \(\mu\) and density f with respect to \(\mu\) . Let q be a lower semi-continuous semi-norm on F.

\begin{enumerate}
\def\labelenumi{\alph{enumi})}
\item
  If, for every compact subset \(\mathrm{K}\) of \(\mathrm{T}\) , the upper integral \(\int_{\mathrm{K}}^{*}(q \circ \mathbf{f}) d\mu\) is finite, then \(\mathbf{m}\) is \(q\) -majorizable.
\item
  If \(\mathbf{m}\) is \(q\) -majorizable, then \(q(\mathbf{m})\) has base \(\mu\) ; if in addition \(\mathbf{f}\) is \(\mu\) -measurable as a mapping of \(\mathrm{T}\) into \(\mathrm{F}_{\sigma}\) , then \(q \circ \mathbf{f}\) is locally \(\mu\) -integrable and \(q(\mathbf{m}) = (q \circ \mathbf{f}) \cdot \mu\) .
\end{enumerate}

\begin{enumerate}
\def\labelenumi{\alph{enumi})}
\tightlist
\item
  For every finite subset \(J\) of \(A_q'\) , let us denote by \(\lambda_J\) the supremum of the measures \(|\mathbf{z}' \circ \mathbf{m}|\) , where \(\mathbf{z}'\) runs over \(J\) ; if \(g_J = \sup_{\mathbf{z}' \in J} |\langle \mathbf{z}', \mathbf{f} \rangle|\) then \(\lambda_{J}=g_{J}\cdot\mu\) (Ch. V, §5, No.~2, Prop. 2). For every relatively compact open subset U of T, let \(\lambda_{J,U}\) be the restriction of \(\lambda_{J}\) to U; let us first show that as J runs over the directed set F of finite subsets of \(A_{q}^{\prime}\) , the family \((\lambda_{J,U})\) is bounded above in \(\mathcal{M}(U)\) . Indeed, for every function \(h\geqslant0\) in \(K(U)\) ,
\end{enumerate}

\[
\int h d \lambda_ {\mathrm{J}, \mathrm{U}} = \int h g _ {\mathrm{J}} d \mu \leqslant \int^ {*} (q \circ \mathbf {f}) h d \mu \leqslant \| h \| \int_ {\mathrm{U}} ^ {*} (q \circ \mathbf {f}) d \mu ,
\]

whence our assertion (Ch. II, §2, No.~2). Let \(\nu_{\mathrm{U}}\) be the supremum of this family of measures in \(\mathcal{M}(\mathrm{U})\) . If \(\mathrm{U}'\) is a second relatively compact open subset of T such that \(\mathrm{U} \subset \mathrm{U}'\) , then \(\nu_{\mathrm{U}}\) is the restriction of \(\nu_{\mathrm{U}'}\) to U, as follows immediately from the expression of the supremum of an increasing directed set of measures (Ch. II, §2, No.~2) and the fact that \(\lambda_{\mathrm{J,U}}\) is the restriction to U of \(\lambda_{\mathrm{J,U}'}\) . Thus there is one and only one positive measure \(\nu\) whose restriction to each U is \(\nu_{\mathrm{U}}\) (Ch. III, §2, No.~1, Prop. 1), and it is clear that \(\nu = q(\mathbf{m})\) .

\begin{enumerate}
\def\labelenumi{\alph{enumi})}
\setcounter{enumi}{1}
\tightlist
\item
  Since the measures \(\lambda_{\mathrm{J}}\) have base \(\mu\) , so does their supremum \(q(\mathbf{m})\) (Ch. V, §5, No.~5, Th. 2). If \(\mathbf{f}\) is \(\mu\) -measurable for the topology \(\sigma(\mathrm{F},\mathrm{F}')\) on \(\mathrm{F}\) , it follows at once from the definitions that the mapping \(g:t\mapsto(g_{\mathrm{J}}(t))_{\mathrm{J}\in\mathfrak{F}}\) of \(\mathrm{T}\) into the product space \(\mathbf{R}^{\mathfrak{F}}\) is \(\mu\) -measurable. The restriction of \(q\circ\mathbf{f}=\sup_{\mathrm{J}\in\mathfrak{F}}g_{\mathrm{J}}\) to every compact subset of \(\mathrm{T}\) on which \(g\) is continuous is lower semi-continuous; consequently \(q \circ f\) is \(\mu\) -measurable (Ch. IV, §5, No.~5, Cor. of Prop. 8 and No.~10, Prop. 15). Let \(K\) be a compact subset of T; it admits a partition consisting of a \(\mu\) -negligible set and a sequence \((\mathrm{K}_n)\) of compact sets on which \(g\) is continuous. Then \(\int_{\mathrm{K}_n}^*\left(q\circ \mathbf{f}\right)d\mu = \sup_{\mathrm{J}}\int_{\mathrm{K}_n}g_{\mathrm{J}}d\mu \leqslant \int_{\mathrm{K}_n}dq(\mathbf{m})\) for all \(n\) (Ch. IV, §1, No.~1, Th. 1 and Ch. V, §7, No.~1, Prop. 1), whence \(\int_{\mathrm{K}}^{*}(q\circ \mathbf{f})d\mu = \sum_{n}\int_{\mathrm{K}_n}^*\left(q\circ \mathbf{f}\right)d\mu \leqslant \int_{\mathrm{K}}dq(\mathbf{m})\) . But this proves that \(q\circ \mathbf{f}\) is locally \(\mu\) -integrable and that \(\lambda_{\mathrm{J}}\leqslant (q\circ \mathbf{f})\cdot \mu \leqslant q(\mathbf{m})\) for all \(\mathbf{J}\in \mathfrak{F}\) ; whence, by definition, \(q(\mathbf{m}) = (q\circ \mathbf{f})\cdot \mu\) .
\end{enumerate}

Remark. --- Suppose that there exists in \(A_{q}^{\prime}\) a countable subset D that is dense for \(\sigma(\mathbf{F}^{\prime},\mathbf{F})\) ; then, the function \(q\circ f\) is always \(\mu\) -measurable, because \(q(\mathbf{f}(t)) = \sup_{\mathbf{z}^{\prime}\in\mathrm{D}} |\langle\mathbf{z}^{\prime},\mathbf{f}(t)\rangle|\) (Ch. IV, §5, No.~4, Cor. 1 of Th. 2). Then, for every

compact subset K of T, \(\int_{K}^{*}(q\circ\mathbf{f})d\mu=\sup_{J}\int_{K}g_{J}d\mu\) , where J runs over the countable directed set of finite subsets of D (Ch. IV, §1, No.~1, Cor. of Th. 3); thus, one sees that in this case the condition that \(\int_{K}^{*}(q\circ\mathbf{f})d\mu<+\infty\) for every compact subset K of T is necessary and sufficient for m to be q-majorizable.

PROPOSITION 9. --- Let F be a finite-dimensional Banach space. Every measure m on T with values in F is a measure with base \textbar m\textbar. If \(m = f \cdot |m|\) , then \(|f(t)| = 1\) locally almost everywhere for \textbar m\textbar. For m to have base \(\mu\) , it is necessary and sufficient that \textbar m\textbar{} have base \(\mu\) , and if \(m = g \cdot \mu\) then \(|m| = |g| \cdot \mu\) .

Let \((\mathbf{e}_i)_{1 \leqslant i \leqslant n}\) and \((\mathbf{e}_i')_{1 \leqslant i \leqslant n}\) be dual bases of F and \(\mathbf{F}'\) (A, II, §2, No.~6) with \(|\mathbf{e}_i'| = 1\) for all \(i\) . Then \(|\mathbf{e}_i' \circ \mathbf{m}| \leqslant |\mathbf{m}|\) for every index \(i\) , therefore (Ch. V, §5, No.~5, Th. 2) \(\mathbf{e}_i' \circ \mathbf{m} = h_i \cdot |\mathbf{m}|\) , where \(h_i\) is bounded and \(|\mathbf{m}|\) -measurable. Setting \(\mathbf{h} = \sum_{i=1}^{n} h_i \cdot \mathbf{e}_i\) , we therefore have \(\mathbf{m} = \mathbf{h} \cdot |\mathbf{m}|\) . If \(\mathbf{m} = \mathbf{f} \cdot |\mathbf{m}|\) , Prop. 8 shows that \(|\mathbf{m}| = |\mathbf{f}| \cdot |\mathbf{m}|\) , whence \(|\mathbf{f}(t)| = 1\) locally almost everywhere for \(|\mathbf{m}|\) (Ch. V, §5, No.~3, Cor. 2 of Prop. 3). The final assertion follows at once from Prop. 8.

Remark. --- If \(z = \sum_{i=1}^{n} z_i e_i\) then \(\psi(z_1, \ldots, z_n) = |z|\) is a positively homogeneous continuous function on \(R to the n\) . Setting \(\mu_i = e_i' \circ m = h_i \cdot |m|\) , by definition (Ch. V, §5, No.~9) \(\psi(\mu_1, \ldots, \mu_n) = \psi(h_1, \ldots, h_n) \cdot |m| = |h| \cdot |m| = |m|\) .

